% !PRINTED 496
\weberpage{496}
%% [page furniture, not body text -- source-confirmed 2026-08-10]
%% Neunzehnter Abschnitt. \hfill \S\ 133.

Da 1, 2, 4 die quadratischen Reste von 7; 1, 3, 4, 5, 9 die
quadratischen Reste von 11 sind, so können diese beiden Summen nach
Bd.~I, \S\ 179 bestimmt werden, und ergeben in den beiden Fällen:
\[
i\sqrt{7} \quad \text{und} \quad -i\sqrt{11},
\]
so daß
\weq{18}{\begin{aligned}
w_\nu &= q^{-\frac{1}{7}} \varrho^{3\nu} + \cdots, & n &= 7, \\
w_\nu &= -q^{-\frac{5}{22}} \varrho^{5\nu} + \cdots, & n &= 11.
\end{aligned}}
Hierdurch läßt sich eine obere Grenze ableiten, bis zu der in
$A_\nu$ die Potenzen von $f(\omega)$ höchstens ansteigen können,
wenn man noch beachtet, daß die $A_\nu$ als ganze rationale (und
symmetrische) Funktionen der $w_\nu$ darstellbar sind.

Für $n=7$ ergibt sich so, da alle Exponenten von $f(\omega)$ durch
24 teilbar sein müssen, daß $A_1$, $A_2$, $A_3$, $A_4$, $A_5$, $A_6$
Konstanten sind, während $A_7$ nur die erste Potenz von
$f(\omega)^{24}$, und zwar mit dem Koeffizienten $-1$ [nach (18)],
enthält.

Für $n=11$ muß, wenn $f(\omega)^{12\lambda}$ die höchste in $A_\nu$
vorkommende Potenz von $f(\omega)$ ist,
\[
\lambda \leqq \frac{5\nu}{11}
\]
sein, und überdies muß $\lambda$ mit $\nu$ zugleich gerade oder
ungerade sein. Daher sind $A_2$, $A_4$ konstant, $A_6$, $A_8$
enthalten die höchste Potenz $f(\omega)^{24}$, $A_{10}$ enthält
auch $f(\omega)^{48}$. Da die $A$ mit ungeradem Index nur ungerade
Potenzen von $f(\omega)^{12}$ enthalten, so ist $A_1 = 0$, $A_3$,
$A_5$ enthalten $f(\omega)^{12}$, $A_7$, $A_9$ enthalten bis
$f(\omega)^{36}$, in $A_{11}$ steigt $f(\omega)$ bis zur 60sten
Potenz an.

Für den Fall $n=7$ wollen wir nun die Konstanten vollständig mit
Hilfe der komplexen Multiplikation bestimmen, und zwar genügt dazu
die Betrachtung der Werte der $w_\nu$ für $\omega = i$.

Für $\omega = i$ wird
\[
u = f(i) = \sqrt[4]{2},
\]
und wenn wir diesen Wert in die Gleichung (2) einführen, und
\[
\sqrt[4]{2}\, v = x
\]
setzen, so ergibt sich:
\weq{19}{x^8 - 4x^7 + 28x^4 - 32x + 16 = 0.}
Die linke Seite dieser Gleichung ist das Quadrat des Ausdruckes:
\weq{20}{x^4 - 2x^3 - 2x^2 - 4x + 4,}
%CONT>>
